%!TEX program = lualatex
% SPDX-License-Identifier: CC-BY-SA-4.0
% ./build-beamer.sh --euler
\documentclass[aspectratio=169,11pt]{beamer}
\usetheme{AtlasGrand}
% SPDX-License-Identifier: MIT
% Copyright (c) 2026 the repository contributors
% Presentation-local visual vocabulary. Geometry stays outside the theme.
\colorlet{EulerStroke}{AtlasInk}
\colorlet{EulerGround}{AtlasPaper}
% Generated by tools/euler_geometry.py; edit the generator.
% Six rigid faces, five hinges; all artwork is vector.
\newcommand{\EulerFoldClosed}{%
\begin{tikzpicture}[line join=round,line cap=round]
\path[use as bounding box] (-3.6,-2.1) rectangle (3.6,4.3);
\fill[AtlasGold!20!EulerGround] (-1.40000,0.25000)--(0.40000,0.85000)--(0.40000,2.55000)--(-1.40000,1.95000)--cycle;
\path[pattern=north east lines,pattern color=AtlasGold,opacity=.12] (-1.40000,0.25000)--(0.40000,0.85000)--(0.40000,2.55000)--(-1.40000,1.95000)--cycle;
\draw[EulerStroke,line width=.75pt] (-1.40000,0.25000)--(0.40000,0.85000)--(0.40000,2.55000)--(-1.40000,1.95000)--cycle;
\node[text=EulerStroke,font=\scriptsize\scshape] at (-0.50000,1.40000) {II};
\fill[AtlasGold!12!EulerGround] (-0.40000,-0.85000)--(1.40000,-0.25000)--(0.40000,0.85000)--(-1.40000,0.25000)--cycle;
\path[pattern=north east lines,pattern color=AtlasGold,opacity=.12] (-0.40000,-0.85000)--(1.40000,-0.25000)--(0.40000,0.85000)--(-1.40000,0.25000)--cycle;
\draw[EulerStroke,line width=.75pt] (-0.40000,-0.85000)--(1.40000,-0.25000)--(0.40000,0.85000)--(-1.40000,0.25000)--cycle;
\node[text=EulerStroke,font=\scriptsize\scshape] at (0.00000,0.00000) {I};
\fill[AtlasGold!32!EulerGround] (1.40000,-0.25000)--(0.40000,0.85000)--(0.40000,2.55000)--(1.40000,1.45000)--cycle;
\path[pattern=north east lines,pattern color=AtlasGold,opacity=.12] (1.40000,-0.25000)--(0.40000,0.85000)--(0.40000,2.55000)--(1.40000,1.45000)--cycle;
\draw[EulerStroke,line width=.75pt] (1.40000,-0.25000)--(0.40000,0.85000)--(0.40000,2.55000)--(1.40000,1.45000)--cycle;
\node[text=EulerStroke,font=\scriptsize\scshape] at (0.90000,1.15000) {III};
\fill[AtlasGold!26!EulerGround] (-0.40000,-0.85000)--(-1.40000,0.25000)--(-1.40000,1.95000)--(-0.40000,0.85000)--cycle;
\path[pattern=north east lines,pattern color=AtlasGold,opacity=.12] (-0.40000,-0.85000)--(-1.40000,0.25000)--(-1.40000,1.95000)--(-0.40000,0.85000)--cycle;
\draw[EulerStroke,line width=.75pt] (-0.40000,-0.85000)--(-1.40000,0.25000)--(-1.40000,1.95000)--(-0.40000,0.85000)--cycle;
\node[text=EulerStroke,font=\scriptsize\scshape] at (-0.90000,0.55000) {V};
\fill[AtlasGold!8!EulerGround] (-1.40000,1.95000)--(0.40000,2.55000)--(1.40000,1.45000)--(-0.40000,0.85000)--cycle;
\path[pattern=north east lines,pattern color=AtlasGold,opacity=.12] (-1.40000,1.95000)--(0.40000,2.55000)--(1.40000,1.45000)--(-0.40000,0.85000)--cycle;
\draw[EulerStroke,line width=.75pt] (-1.40000,1.95000)--(0.40000,2.55000)--(1.40000,1.45000)--(-0.40000,0.85000)--cycle;
\node[text=EulerStroke,font=\scriptsize\scshape] at (-0.00000,1.70000) {VI};
\fill[AtlasGold!16!EulerGround] (-0.40000,-0.85000)--(1.40000,-0.25000)--(1.40000,1.45000)--(-0.40000,0.85000)--cycle;
\path[pattern=north east lines,pattern color=AtlasGold,opacity=.12] (-0.40000,-0.85000)--(1.40000,-0.25000)--(1.40000,1.45000)--(-0.40000,0.85000)--cycle;
\draw[EulerStroke,line width=.75pt] (-0.40000,-0.85000)--(1.40000,-0.25000)--(1.40000,1.45000)--(-0.40000,0.85000)--cycle;
\node[text=EulerStroke,font=\scriptsize\scshape] at (0.50000,0.30000) {IV};
\end{tikzpicture}}

\newcommand{\EulerFoldOpening}{%
\begin{tikzpicture}[line join=round,line cap=round]
\path[use as bounding box] (-3.6,-2.1) rectangle (3.6,4.3);
\fill[AtlasGold!20!EulerGround] (-1.40000,0.25000)--(0.40000,0.85000)--(-0.17358,2.87349)--(-1.97358,2.27349)--cycle;
\path[pattern=north east lines,pattern color=AtlasGold,opacity=.12] (-1.40000,0.25000)--(0.40000,0.85000)--(-0.17358,2.87349)--(-1.97358,2.27349)--cycle;
\draw[EulerStroke,line width=.75pt] (-1.40000,0.25000)--(0.40000,0.85000)--(-0.17358,2.87349)--(-1.97358,2.27349)--cycle;
\node[text=EulerStroke,font=\scriptsize\scshape] at (-0.78679,1.56175) {II};
\fill[AtlasGold!32!EulerGround] (1.40000,-0.25000)--(0.40000,0.85000)--(1.43244,2.58670)--(2.43244,1.48670)--cycle;
\path[pattern=north east lines,pattern color=AtlasGold,opacity=.12] (1.40000,-0.25000)--(0.40000,0.85000)--(1.43244,2.58670)--(2.43244,1.48670)--cycle;
\draw[EulerStroke,line width=.75pt] (1.40000,-0.25000)--(0.40000,0.85000)--(1.43244,2.58670)--(2.43244,1.48670)--cycle;
\node[text=EulerStroke,font=\scriptsize\scshape] at (1.41622,1.16835) {III};
\fill[AtlasGold!12!EulerGround] (-0.40000,-0.85000)--(1.40000,-0.25000)--(0.40000,0.85000)--(-1.40000,0.25000)--cycle;
\path[pattern=north east lines,pattern color=AtlasGold,opacity=.12] (-0.40000,-0.85000)--(1.40000,-0.25000)--(0.40000,0.85000)--(-1.40000,0.25000)--cycle;
\draw[EulerStroke,line width=.75pt] (-0.40000,-0.85000)--(1.40000,-0.25000)--(0.40000,0.85000)--(-1.40000,0.25000)--cycle;
\node[text=EulerStroke,font=\scriptsize\scshape] at (0.00000,0.00000) {I};
\fill[AtlasGold!8!EulerGround] (-1.97358,2.27349)--(-0.17358,2.87349)--(0.16844,4.09475)--(-1.63156,3.49475)--cycle;
\path[pattern=north east lines,pattern color=AtlasGold,opacity=.12] (-1.97358,2.27349)--(-0.17358,2.87349)--(0.16844,4.09475)--(-1.63156,3.49475)--cycle;
\draw[EulerStroke,line width=.75pt] (-1.97358,2.27349)--(-0.17358,2.87349)--(0.16844,4.09475)--(-1.63156,3.49475)--cycle;
\node[text=EulerStroke,font=\scriptsize\scshape] at (-0.90257,3.18412) {VI};
\fill[AtlasGold!26!EulerGround] (-0.40000,-0.85000)--(-1.40000,0.25000)--(-2.43244,1.29841)--(-1.43244,0.19841)--cycle;
\path[pattern=north east lines,pattern color=AtlasGold,opacity=.12] (-0.40000,-0.85000)--(-1.40000,0.25000)--(-2.43244,1.29841)--(-1.43244,0.19841)--cycle;
\draw[EulerStroke,line width=.75pt] (-0.40000,-0.85000)--(-1.40000,0.25000)--(-2.43244,1.29841)--(-1.43244,0.19841)--cycle;
\node[text=EulerStroke,font=\scriptsize\scshape] at (-1.41622,0.22421) {V};
\fill[AtlasGold!16!EulerGround] (-0.40000,-0.85000)--(1.40000,-0.25000)--(1.97358,0.51162)--(0.17358,-0.08838)--cycle;
\path[pattern=north east lines,pattern color=AtlasGold,opacity=.12] (-0.40000,-0.85000)--(1.40000,-0.25000)--(1.97358,0.51162)--(0.17358,-0.08838)--cycle;
\draw[EulerStroke,line width=.75pt] (-0.40000,-0.85000)--(1.40000,-0.25000)--(1.97358,0.51162)--(0.17358,-0.08838)--cycle;
\node[text=EulerStroke,font=\scriptsize\scshape] at (0.78679,-0.16919) {IV};
\end{tikzpicture}}

\newcommand{\EulerFoldOpen}{%
\begin{tikzpicture}[line join=round,line cap=round]
\path[use as bounding box] (-3.6,-2.1) rectangle (3.6,4.3);
\fill[AtlasGold!8!EulerGround] (-2.40000,1.35000)--(-0.60000,1.95000)--(-1.60000,3.05000)--(-3.40000,2.45000)--cycle;
\path[pattern=north east lines,pattern color=AtlasGold,opacity=.12] (-2.40000,1.35000)--(-0.60000,1.95000)--(-1.60000,3.05000)--(-3.40000,2.45000)--cycle;
\draw[EulerStroke,line width=.75pt] (-2.40000,1.35000)--(-0.60000,1.95000)--(-1.60000,3.05000)--(-3.40000,2.45000)--cycle;
\node[text=EulerStroke,font=\scriptsize\scshape] at (-2.00000,2.20000) {VI};
\fill[AtlasGold!20!EulerGround] (-1.40000,0.25000)--(0.40000,0.85000)--(-0.60000,1.95000)--(-2.40000,1.35000)--cycle;
\path[pattern=north east lines,pattern color=AtlasGold,opacity=.12] (-1.40000,0.25000)--(0.40000,0.85000)--(-0.60000,1.95000)--(-2.40000,1.35000)--cycle;
\draw[EulerStroke,line width=.75pt] (-1.40000,0.25000)--(0.40000,0.85000)--(-0.60000,1.95000)--(-2.40000,1.35000)--cycle;
\node[text=EulerStroke,font=\scriptsize\scshape] at (-1.00000,1.10000) {II};
\fill[AtlasGold!32!EulerGround] (1.40000,-0.25000)--(0.40000,0.85000)--(2.20000,1.45000)--(3.20000,0.35000)--cycle;
\path[pattern=north east lines,pattern color=AtlasGold,opacity=.12] (1.40000,-0.25000)--(0.40000,0.85000)--(2.20000,1.45000)--(3.20000,0.35000)--cycle;
\draw[EulerStroke,line width=.75pt] (1.40000,-0.25000)--(0.40000,0.85000)--(2.20000,1.45000)--(3.20000,0.35000)--cycle;
\node[text=EulerStroke,font=\scriptsize\scshape] at (1.80000,0.60000) {III};
\fill[AtlasGold!12!EulerGround] (-0.40000,-0.85000)--(1.40000,-0.25000)--(0.40000,0.85000)--(-1.40000,0.25000)--cycle;
\path[pattern=north east lines,pattern color=AtlasGold,opacity=.12] (-0.40000,-0.85000)--(1.40000,-0.25000)--(0.40000,0.85000)--(-1.40000,0.25000)--cycle;
\draw[EulerStroke,line width=.75pt] (-0.40000,-0.85000)--(1.40000,-0.25000)--(0.40000,0.85000)--(-1.40000,0.25000)--cycle;
\node[text=EulerStroke,font=\scriptsize\scshape] at (0.00000,0.00000) {I};
\fill[AtlasGold!26!EulerGround] (-0.40000,-0.85000)--(-1.40000,0.25000)--(-3.20000,-0.35000)--(-2.20000,-1.45000)--cycle;
\path[pattern=north east lines,pattern color=AtlasGold,opacity=.12] (-0.40000,-0.85000)--(-1.40000,0.25000)--(-3.20000,-0.35000)--(-2.20000,-1.45000)--cycle;
\draw[EulerStroke,line width=.75pt] (-0.40000,-0.85000)--(-1.40000,0.25000)--(-3.20000,-0.35000)--(-2.20000,-1.45000)--cycle;
\node[text=EulerStroke,font=\scriptsize\scshape] at (-1.80000,-0.60000) {V};
\fill[AtlasGold!16!EulerGround] (-0.40000,-0.85000)--(1.40000,-0.25000)--(2.40000,-1.35000)--(0.60000,-1.95000)--cycle;
\path[pattern=north east lines,pattern color=AtlasGold,opacity=.12] (-0.40000,-0.85000)--(1.40000,-0.25000)--(2.40000,-1.35000)--(0.60000,-1.95000)--cycle;
\draw[EulerStroke,line width=.75pt] (-0.40000,-0.85000)--(1.40000,-0.25000)--(2.40000,-1.35000)--(0.60000,-1.95000)--cycle;
\node[text=EulerStroke,font=\scriptsize\scshape] at (1.00000,-1.10000) {IV};
\end{tikzpicture}}


% Sparse pages deliberately suspend the grand border and paper texture.
\newcommand{\EulerQuiet}{%
  \setbeamertemplate{background canvas}{\color{AtlasPaper}\vrule width\paperwidth height\paperheight}%
  \setbeamertemplate{footline}{}%
}
\newcommand{\EulerNight}{%
  \AtlasGrandNight
  \colorlet{EulerStroke}{AtlasGoldLight}\colorlet{EulerGround}{AtlasNight}}
\newcommand{\EulerCaption}[2]{%
  {\small\color{AtlasGold}\textls[140]{\textsc{#1}}\par}\smallskip
  {\large\itshape #2\par}}

% The cube graph in a Schlegel drawing. The five removed edges are a
% complement of the path a-b-c-d-h-g-f-e. Every prefix remains connected.
\newcommand{\EulerMap}[2]{% deleted cycle edges (0..5), pruned leaves (0..7)
  \begin{tikzpicture}[line cap=round,line join=round]
    \path[use as bounding box] (-3.1,-3.05) rectangle (3.1,3.05);
    \foreach \name/\x/\y in {a/-2.35/2.35,b/2.35/2.35,c/2.35/-2.35,d/-2.35/-2.35,e/-.9/.9,f/.9/.9,g/.9/-.9,h/-.9/-.9}
      \coordinate (\name) at (\x,\y);
    % Face tint disappears when that region joins the exterior. The
    % innermost face is opened last; all five face mergers are explicit.
    \ifnum#1<1\fill[AtlasGold!12!EulerGround] (a)--(d)--(h)--(e)--cycle;\fi
    \ifnum#1<2\fill[AtlasGold!19!EulerGround] (a)--(b)--(f)--(e)--cycle;\fi
    \ifnum#1<3\fill[AtlasGold!27!EulerGround] (b)--(c)--(g)--(f)--cycle;\fi
    \ifnum#1<4\fill[AtlasGold!35!EulerGround] (c)--(d)--(h)--(g)--cycle;\fi
    \ifnum#1<5\fill[AtlasRubric!12!EulerGround] (e)--(f)--(g)--(h)--cycle;\fi
    \foreach \u/\v/\step in {d/a/1,a/e/2,b/f/3,c/g/4,e/h/5}{
      \ifnum#1<\step
        \draw[EulerStroke,line width=1pt] (\u)--(\v);
      \fi
    }
    \foreach \u/\v/\step in {a/b/1,b/c/2,c/d/3,d/h/4,h/g/5,g/f/6,f/e/7}{
      \ifnum#2<\step
        \draw[EulerStroke,line width=1.2pt] (\u)--(\v);
        \draw[AtlasGold,line width=.25pt] (\u)--(\v);
      \fi
    }
    \foreach \n/\step in {a/1,b/2,c/3,d/4,h/5,g/6,f/7,e/8}{
      \ifnum#2<\step
        \fill[EulerGround] (\n) circle (.10);
        \draw[EulerStroke,line width=.5pt,fill=AtlasGold] (\n) circle (.065);
      \fi
    }
    \ifnum#1=0
      \node[font=\scriptsize\itshape,text=AtlasRubric] at (0,0) {one face};
      \node[font=\scriptsize\itshape,text=AtlasRubric] at (0,2.8) {the exterior is a face, too};
    \fi
  \end{tikzpicture}%
}

\newcommand{\EulerLedger}[3]{%
  \begin{tabular}{@{}ccc@{}}
    {\color{AtlasGold}\scshape vertices}&{\color{AtlasGold}\scshape edges}&{\color{AtlasGold}\scshape faces}\\[1mm]
    {\Large #1}&{\Large #2}&{\Large #3}
  \end{tabular}\par
}
\newcommand{\EulerCount}[3]{%
  \EulerLedger{#1}{#2}{#3}\medskip
  {\large $#1-#2+#3$\par}}

% A purpose-built solid, used for the cover and the counting plate.
\newcommand{\EulerCube}[1]{% 0 engraved, 1 vertices, 2 edges, 3 faces
  \begin{tikzpicture}[line join=round,line cap=round]
    \path[use as bounding box] (-2.7,-2.4) rectangle (2.7,3.3);
    \foreach \n/\x/\y in {a/-2/-.7,b/.6/-1.55,c/2/-.3,d/-.6/.55,e/-2/1.8,f/.6/.95,g/2/2.2,h/-.6/3.05}
      \coordinate (\n) at (\x,\y);
    \fill[AtlasGold!18!EulerGround] (a)--(b)--(f)--(e)--cycle;
    \fill[AtlasGold!34!EulerGround] (b)--(c)--(g)--(f)--cycle;
    \fill[AtlasGold!8!EulerGround] (e)--(f)--(g)--(h)--cycle;
    \begin{scope}
      \clip (a)--(b)--(c)--(g)--(h)--(e)--cycle;
      \foreach \y in {-3,-2.88,...,4}
        \draw[AtlasGold,opacity=.23,line width=.22pt] (-3,\y)--++(6,1.3);
    \end{scope}
    \draw[EulerStroke,opacity=.45,densely dashed,line width=.5pt] (a)--(d)--(c) (d)--(h);
    \ifnum#1=2\def\ew{1.8pt}\else\def\ew{.8pt}\fi
    \draw[EulerStroke,line width=\ew] (a)--(b)--(c)--(g)--(h)--(e)--cycle
      (b)--(f)--(e) (f)--(g);
    \ifnum#1=1
      \foreach \n in {a,b,c,d,e,f,g,h}{
        \fill[EulerGround] (\n) circle (.14);
        \draw[EulerStroke,fill=AtlasGold,line width=.5pt] (\n) circle (.085);
      }
    \fi
    \ifnum#1=3
      \foreach \x/\y in {-.7/.15,1.3/.4,0/2}
        \node[text=EulerStroke,font=\small\itshape] at (\x,\y) {face};
    \fi
  \end{tikzpicture}%
}

\newcommand{\EulerSphere}{%
  \begin{tikzpicture}[draw=EulerStroke,line width=.45pt]
    \path[use as bounding box] (-2.6,-2.6) rectangle (2.6,2.6);
    \foreach \a in {-60,-30,0,30,60}{
      \pgfmathsetmacro{\yy}{2.3*sin(\a)}
      \pgfmathsetmacro{\xx}{2.3*cos(\a)}
      \draw[opacity=.4] (0,\yy) ellipse (\xx cm and .18cm);
    }
    \foreach \r in {.55,1.15,1.8}\draw[opacity=.45] (0,0) ellipse (\r cm and 2.3cm);
    \draw[line width=1pt] (0,0) circle (2.3);
    \draw (0,0) ellipse (2.3cm and .5cm);
  \end{tikzpicture}%
}

\title{The Form of Two}
\subtitle{A polyhedron, unfolded into an invariant}
\author{latex-atlas}
\date{November 2026}

\begin{document}

% I. An object before a formula. The cover uses scale, not a title box.
{\EulerNight
\begin{frame}[plain,noframenumbering]
  \begin{tikzpicture}[remember picture,overlay]
    \node[anchor=north west,text=AtlasGoldLight,font=\small]
      at ([xshift=15mm,yshift=-14mm]current page.north west)
      {\textls[150]{\textsc{A geometry of what remains}}};
    \node[anchor=north west,text=AtlasIvory,align=left,inner sep=0pt]
      at ([xshift=16mm,yshift=-29mm]current page.north west)
      {{\fontsize{38}{39}\selectfont\scshape The Form}};
    \node[anchor=north west,text=AtlasIvory,inner sep=0pt]
      at ([xshift=16mm,yshift=-45mm]current page.north west)
      {{\fontsize{38}{39}\selectfont\scshape of Two}};
    \node[anchor=north west,text=AtlasGoldLight,text width=64mm,align=left,
      font=\large\itshape,inner sep=0pt]
      at ([xshift=17mm,yshift=-64mm]current page.north west)
      {A polyhedron, unfolded\\into an invariant.};
    \node at ([xshift=118mm,yshift=-45mm]current page.north west)
      {\resizebox{!}{68mm}{\EulerCube{0}}};
    \node[anchor=south west,text=AtlasGoldLight,font=\scriptsize]
      at ([xshift=16mm,yshift=9mm]current page.south west)
      {\textls[150]{\textsc{Atlas\quad /\quad Geometria\quad /\quad MMXXVI}}};
  \end{tikzpicture}
\end{frame}}

\section{The object}
\begin{frame}{Three ways of seeing}
  \begin{columns}[c,onlytextwidth]
    \begin{column}{.53\textwidth}
      \centering
      \only<1|handout:0>{\resizebox{!}{53mm}{\EulerCube{1}}}%
      \only<2|handout:0>{\resizebox{!}{53mm}{\EulerCube{2}}}%
      \only<3|handout:1>{\resizebox{!}{53mm}{\EulerCube{3}}}%
    \end{column}
    \begin{column}{.41\textwidth}
      \begin{overlayarea}{\linewidth}{49mm}
        \AtlasGrandKicker{An inventory of the cube}\medskip
        \uncover<1->{\noindent{\fontsize{29}{32}\selectfont $8$}\quad vertices\par\medskip}
        \uncover<2->{\noindent{\fontsize{29}{32}\selectfont $12$}\quad edges\par\medskip}
        \uncover<3->{\noindent{\fontsize{29}{32}\selectfont $6$}\quad faces\par\medskip
          {\small\itshape Including the parts hidden from view.}}
      \end{overlayarea}
    \end{column}
  \end{columns}
\end{frame}

\begin{frame}{Let the surface open}
  \framesubtitle{Six rigid squares; five surviving hinges}
  \begin{columns}[c,onlytextwidth]
    \begin{column}{.56\textwidth}
      \centering
      \only<1|handout:0>{\resizebox{!}{51mm}{\EulerFoldClosed}}%
      \only<2|handout:0>{\resizebox{!}{51mm}{\EulerFoldOpening}}%
      \only<3|handout:1>{\resizebox{!}{51mm}{\EulerFoldOpen}}%
    \end{column}
    \begin{column}{.38\textwidth}
      \AtlasGrandKicker{The cut-out net}
      The faces turn about their hinges. Lengths and angles within each
      square survive.\par\medskip
      \AtlasGrandScholium{Keep the identities}{Cut edges appear twice in
        the net. This is a view of the surface, not the graph on which
        we will count.}
    \end{column}
  \end{columns}
\end{frame}

\begin{frame}{Give one face to the horizon}
  \framesubtitle{The cube redrawn in the plane, without cutting its edges}
  \begin{columns}[c,onlytextwidth]
    \begin{column}{.52\textwidth}
      \centering\resizebox{!}{54mm}{\EulerMap{0}{0}}
    \end{column}
    \begin{column}{.42\textwidth}
      \begin{overlayarea}{\linewidth}{48mm}
        \AtlasGrandKicker{A change of viewpoint}
        The four corners of one face become the outer square. The
        opposite face sits inside it.\par\medskip
        \uncover<2->{The chosen face becomes the unbounded exterior.
          \emph{It still counts.}\par\medskip
          \EulerLedger{8}{12}{6}}
      \end{overlayarea}
    \end{column}
  \end{columns}
\end{frame}

% V. A deliberate silence between the object and the operation.
{\EulerQuiet
\begin{frame}[plain]
  \centering
  \EulerCaption{The first operation}{Erase an edge on a cycle.}\bigskip
  {\fontsize{34}{40}\selectfont $\Delta E=-1\qquad\Delta F=-1$\par}
  \bigskip
  \begin{minipage}{103mm}\centering
    Its two incident faces become one. The rest of the cycle keeps
    the graph connected. No vertex is lost.
  \end{minipage}\par\bigskip
  {\color{AtlasRubric}\Large $V-E+F\quad\text{is unchanged.}$\par}
\end{frame}}

\section{The transformation}
\begin{frame}{One boundary less; one region less}
  \begin{columns}[c,onlytextwidth]
    \begin{column}{.53\textwidth}
      \centering
      \only<1|handout:0>{\resizebox{!}{54mm}{\EulerMap{0}{0}}}%
      \only<2|handout:0>{\resizebox{!}{54mm}{\EulerMap{1}{0}}}%
      \only<3|handout:0>{\resizebox{!}{54mm}{\EulerMap{2}{0}}}%
      \only<4|handout:0>{\resizebox{!}{54mm}{\EulerMap{3}{0}}}%
      \only<5|handout:0>{\resizebox{!}{54mm}{\EulerMap{4}{0}}}%
      \only<6|handout:1>{\resizebox{!}{54mm}{\EulerMap{5}{0}}}%
    \end{column}
    \begin{column}{.41\textwidth}
      \begin{overlayarea}{\linewidth}{54mm}
        \AtlasGrandKicker{The regions join the exterior}\medskip
        \only<1|handout:0>{\EulerCount{8}{12}{6}}%
        \only<2|handout:0>{\EulerCount{8}{11}{5}}%
        \only<3|handout:0>{\EulerCount{8}{10}{4}}%
        \only<4|handout:0>{\EulerCount{8}{9}{3}}%
        \only<5|handout:0>{\EulerCount{8}{8}{2}}%
        \only<6|handout:1>{\EulerCount{8}{7}{1}}%
        \medskip
        {\small\raggedright Repeat while a cycle remains. Each erased edge
          merges two faces; the connected drawing eventually becomes a tree.\par}
        \only<beamer:0|handout:1>{\smallskip{\scriptsize
          $(E,F): (12,6)\to(11,5)\to(10,4)$\par
          $\phantom{(E,F):}\to(9,3)\to(8,2)\to(7,1)$.}}
      \end{overlayarea}
    \end{column}
  \end{columns}
\end{frame}

% VII. Full scale, with all ornament withheld. A tree has no enclosed face.
{\EulerQuiet
\begin{frame}[plain]
  \begin{tikzpicture}[remember picture,overlay]
    \node at ([xshift=52mm,yshift=-43mm]current page.north west)
      {\resizebox{!}{69mm}{\EulerMap{5}{0}}};
    \node[anchor=west,text width=53mm,align=left,text=AtlasInk]
      at ([xshift=98mm,yshift=-40mm]current page.north west)
      {{\fontsize{27}{29}\selectfont\itshape Nothing\\is enclosed.}\par\bigskip
       {\large One face remains:\\the whole exterior.}\par\bigskip
       {\color{AtlasRubric}\Large $F=1$}};
    \node[anchor=south west,text=AtlasGold,font=\scriptsize]
      at ([xshift=16mm,yshift=8mm]current page.south west)
      {\textls[160]{\textsc{The tree beneath the surface}}};
  \end{tikzpicture}
\end{frame}}

\begin{frame}{Now let the branches fall}
  \framesubtitle{The second operation: remove a leaf and its incident edge}
  \begin{columns}[c,onlytextwidth]
    \begin{column}{.53\textwidth}
      \centering
      \only<1|handout:0>{\resizebox{!}{51mm}{\EulerMap{5}{0}}}%
      \only<2|handout:0>{\resizebox{!}{51mm}{\EulerMap{5}{1}}}%
      \only<3|handout:0>{\resizebox{!}{51mm}{\EulerMap{5}{2}}}%
      \only<4|handout:0>{\resizebox{!}{51mm}{\EulerMap{5}{3}}}%
      \only<5|handout:0>{\resizebox{!}{51mm}{\EulerMap{5}{4}}}%
      \only<6|handout:0>{\resizebox{!}{51mm}{\EulerMap{5}{5}}}%
      \only<7|handout:0>{\resizebox{!}{51mm}{\EulerMap{5}{6}}}%
      \only<8|handout:1>{\resizebox{!}{51mm}{\EulerMap{5}{7}}}%
    \end{column}
    \begin{column}{.41\textwidth}
      \begin{overlayarea}{\linewidth}{51mm}
        \AtlasGrandKicker{Vertices and edges depart together}\smallskip
        \only<1|handout:0>{\EulerCount{8}{7}{1}}%
        \only<2|handout:0>{\EulerCount{7}{6}{1}}%
        \only<3|handout:0>{\EulerCount{6}{5}{1}}%
        \only<4|handout:0>{\EulerCount{5}{4}{1}}%
        \only<5|handout:0>{\EulerCount{4}{3}{1}}%
        \only<6|handout:0>{\EulerCount{3}{2}{1}}%
        \only<7|handout:0>{\EulerCount{2}{1}{1}}%
        \only<8|handout:1>{\EulerCount{1}{0}{1}}%
        \smallskip
        {\small\raggedright Every finite tree with at least two vertices
          has a leaf. Removing it leaves a smaller tree, with the same
          single face.\par}\smallskip
        {\color{AtlasRubric}$\Delta V=\Delta E=-1,\quad\Delta F=0.$}
      \end{overlayarea}
    \end{column}
  \end{columns}
\end{frame}

{\EulerQuiet
\begin{frame}[plain]
  \centering
  \EulerCaption{The irreducible drawing}{One point. No edges. One surrounding face.}
  \bigskip
  \begin{tikzpicture}
    \path[use as bounding box] (-3,-.8) rectangle (3,.8);
    \draw[AtlasGold,line width=.4pt] (0,0) circle (.22);
    \fill[AtlasInk] (0,0) circle (.07);
  \end{tikzpicture}\par\bigskip
  {\fontsize{34}{40}\selectfont $1-0+1$\par}
  \bigskip
  {\large\itshape The expression has kept its value through every loss.\par}
\end{frame}}

% X. The answer gets the largest type and the full illuminated field.
{\EulerNight
\begin{frame}[plain]
  \begin{tikzpicture}[remember picture,overlay]
    \node[text=AtlasGoldLight,font=\small] at
      ([yshift=25mm]current page.center)
      {\textls[200]{\textsc{Euler's polyhedral formula}}};
    \node[text=AtlasIvory,inner sep=0pt] at
      ([yshift=2mm]current page.center)
      {{\fontsize{58}{65}\selectfont $V-E+F=2$}};
    \node[text=AtlasGoldLight,font=\large\itshape] at
      ([yshift=-23mm]current page.center)
      {What changed was the drawing. What endured was the number.};
    \node at ([yshift=-33mm]current page.center)
      {\pgfornament[width=25mm,color=AtlasGoldLight]{71}};
  \end{tikzpicture}
\end{frame}}

\section{The invariant}
\begin{frame}{Many forms; the same remainder}
  \begin{columns}[c,onlytextwidth]
    \begin{column}{.30\textwidth}\centering
      \begin{tikzpicture}[scale=.85,draw=AtlasInk,line width=.7pt]
        \path[use as bounding box] (-1.9,-1.8) rectangle (1.9,2);
        \fill[AtlasGold!16!AtlasPaper] (0,1.7)--(-1.6,-1)--(.5,-1.5)--cycle;
        \fill[AtlasGold!30!AtlasPaper] (0,1.7)--(.5,-1.5)--(1.5,-.5)--cycle;
        \draw[dashed,opacity=.4] (-1.6,-1)--(1.5,-.5);
        \draw (0,1.7)--(-1.6,-1)--(.5,-1.5)--(1.5,-.5)--cycle (0,1.7)--(.5,-1.5);
      \end{tikzpicture}\par
      \AtlasGrandKicker{Tetrahedron}\medskip
      {\Large $4-6+4=2$}
    \end{column}
    \begin{column}{.30\textwidth}\centering
      \resizebox{!}{32.3mm}{\EulerCube{0}}\par
      \AtlasGrandKicker{Cube}\medskip
      {\Large $8-12+6=2$}
    \end{column}
    \begin{column}{.30\textwidth}\centering
      \begin{tikzpicture}[scale=.85,draw=AtlasInk,line width=.7pt]
        \path[use as bounding box] (-1.9,-1.8) rectangle (1.9,2);
        \fill[AtlasGold!16!AtlasPaper] (0,1.7)--(-1.6,0)--(.4,-.6)--cycle;
        \fill[AtlasGold!30!AtlasPaper] (0,1.7)--(.4,-.6)--(1.6,0)--cycle;
        \fill[AtlasGold!24!AtlasPaper] (0,-1.6)--(-1.6,0)--(.4,-.6)--cycle;
        \fill[AtlasGold!40!AtlasPaper] (0,-1.6)--(.4,-.6)--(1.6,0)--cycle;
        \draw[dashed,opacity=.4] (-1.6,0)--(-.4,.5)--(1.6,0) (0,1.7)--(-.4,.5)--(0,-1.6);
        \draw (0,1.7)--(-1.6,0)--(0,-1.6)--(1.6,0)--cycle
          (-1.6,0)--(.4,-.6)--(1.6,0) (0,1.7)--(.4,-.6)--(0,-1.6);
      \end{tikzpicture}\par
      \AtlasGrandKicker{Octahedron}\medskip
      {\Large $6-12+8=2$}
    \end{column}
  \end{columns}
  \bigskip\bigskip
  \centerline{\itshape The cube illustrated the proof. The operations did not depend on the cube.}
\end{frame}

{\EulerNight
\begin{frame}[plain]
  \begin{tikzpicture}[remember picture,overlay]
    \node at ([xshift=43mm,yshift=-43mm]current page.north west)
      {\resizebox{!}{55mm}{\EulerSphere}};
    \node[anchor=north west,text=AtlasGoldLight,font=\small] at
      ([xshift=80mm,yshift=-18mm]current page.north west)
      {\textls[150]{\textsc{The shape of the theorem}}};
    \node[anchor=north west,text=AtlasIvory,text width=65mm,align=left,inner sep=0pt] at
      ([xshift=81mm,yshift=-28mm]current page.north west)
      {{\fontsize{20}{24}\selectfont\scshape Every finite\\connected plane graph.}\par\medskip
       Count the exterior among its faces. Then $V-E+F=2$.\par\medskip
       {\small A convex polyhedron has a spherical surface, so its edges
       admit just such a drawing. A surface with handles is another story.}};
    \node[anchor=south west,text=AtlasGoldLight,font=\scriptsize] at
      ([xshift=16mm,yshift=9mm]current page.south west)
      {\textls[100]{\textsc{Cycles\enspace $\longrightarrow$\enspace tree\enspace $\longrightarrow$\enspace point}}};
    \node[anchor=south east,text=AtlasGoldLight,font=\scriptsize] at
      ([xshift=-14mm,yshift=9mm]current page.south east)
      {\href{https://ics.uci.edu/~eppstein/junkyard/euler/iface.html}{\textsc{Proof notes: Eppstein's Geometry Junkyard}}};
  \end{tikzpicture}
\end{frame}}

\end{document}
